\section{Fixed population or all flights}\label{sec:fixed}

Does letting the population shrink with the lag change the answer?

In design B the flights that enter a lag are those with a piece of at least $\tau/0.8$ s. They are not a random subset: they are
the flights with the longest continuous pieces, and their MSD is larger. Along a curve the MSD of design B thus mixes the growth of
the MSD of one flight with a change of the flights, and $\alpha(\tau)$ bends upward where the population is small
(Fig.~\ref{fig:fixed-vs-all}). On the same lags, $H$ of design B is below or equal to that of design A in every population and both
disciplines (equal in hills, where $T=$ \UniTauFixed{} s) (Table~\ref{tab:fixed-vs-all-H}); for all terrain and the lags up to \UniTAllAlongest{} s, \UniHFirstAllAlongest{} against
\UniHFirstAllBlongest{} for paragliders and \UniHSecondAllAlongest{} against \UniHSecondAllBlongest{} for hang gliders.

\begin{figure}
\centering
\includegraphics[width=\textwidth]{fixed_vs_all.pdf}
\caption{A shrinking population bends $\alpha(\tau)$ upward at long lags, which the fixed population does not do. All terrain,
lags up to the longest $T$; solid: fixed population (A); dash-dot: all flights (B); the dashed and dotted horizontal lines are
$2H$ of A and of B. Both panels share the vertical axis, clipped at 2.8 (the hang band exceeds it above $10^4$ s).}\label{fig:fixed-vs-all}
\end{figure}

\begin{table}
\centering
\small
\renewcommand{\arraystretch}{1.15}\setlength{\tabcolsep}{4pt}
\begin{adjustbox}{max width=\textwidth}
\begin{tabular}{lrrrrr}
\toprule
Population & $T$ (s) & para, A & para, B & hang, A & hang, B \\
\midrule
all & 11450 & 0.923 & 0.917 & 0.915 & 0.897 \\
plains & 4150 & 0.946 & 0.944 & 0.933 & 0.931 \\
hills & 1000 & 0.927 & 0.927 & 0.909 & 0.909 \\
low mountains & 9350 & 0.902 & 0.893 & 0.861 & 0.851 \\
high mountains & 7630 & 0.897 & 0.889 & 0.869 & 0.861 \\
Alps & 9350 & 0.895 & 0.887 & 0.866 & 0.855 \\
\bottomrule
\end{tabular}
\end{adjustbox}
\caption{$H$ of each discipline over the lags up to the longest $T$ of the population: fixed population (A) and all flights (B).}\label{tab:fixed-vs-all-H}
\end{table}


The question is whether the difference between the disciplines depends on the design. Up to $1000$ s, where the two designs
differ only by the flights whose pieces are shorter than \UniPieceFixed{} s, the differences $\Hpara-\Hhang$ of A and B agree to within
\UniGapKilo{} in every population (Table~\ref{tab:fixed-vs-all}, Fig.~\ref{fig:forest}). Up to the longest $T$ the point values differ by up to \UniGapLongest, \UniDHAllBlongest{} against \UniDHAllAlongest{} for all
terrain, since the hang glider flights fall to \UniNSecondAllBlongest; the bands of A and B overlap in every population.

\begin{table}
\centering
\small
\renewcommand{\arraystretch}{1.15}\setlength{\tabcolsep}{4pt}
\begin{adjustbox}{max width=\textwidth}
\begin{tabular}{lcccc}
\toprule
Population & A, $T=1000$ s & A, longest $T$ & B, to 1000 s & B, to longest $T$ \\
\midrule
all & \shortstack{$+$0.024\\[-1pt]{\scriptsize [$+$0.017, $+$0.031]}} & \shortstack{$+$0.007\\[-1pt]{\scriptsize [$-$0.006, $+$0.020]}} & \shortstack{$+$0.024\\[-1pt]{\scriptsize [$+$0.017, $+$0.032]}} & \shortstack{$+$0.020\\[-1pt]{\scriptsize [$+$0.006, $+$0.033]}} \\
plains & \shortstack{$+$0.009\\[-1pt]{\scriptsize [$+$0.000, $+$0.018]}} & \shortstack{$+$0.013\\[-1pt]{\scriptsize [$+$0.005, $+$0.021]}} & \shortstack{$+$0.010\\[-1pt]{\scriptsize [$+$0.000, $+$0.020]}} & \shortstack{$+$0.013\\[-1pt]{\scriptsize [$+$0.005, $+$0.021]}} \\
hills & \shortstack{$+$0.018\\[-1pt]{\scriptsize [$+$0.010, $+$0.025]}} & \shortstack{$+$0.018\\[-1pt]{\scriptsize [$+$0.010, $+$0.025]}} & \shortstack{$+$0.017\\[-1pt]{\scriptsize [$+$0.010, $+$0.025]}} & \shortstack{$+$0.017\\[-1pt]{\scriptsize [$+$0.010, $+$0.025]}} \\
low mountains & \shortstack{$+$0.027\\[-1pt]{\scriptsize [$+$0.022, $+$0.032]}} & \shortstack{$+$0.040\\[-1pt]{\scriptsize [$+$0.030, $+$0.050]}} & \shortstack{$+$0.027\\[-1pt]{\scriptsize [$+$0.022, $+$0.032]}} & \shortstack{$+$0.042\\[-1pt]{\scriptsize [$+$0.035, $+$0.049]}} \\
high mountains & \shortstack{$+$0.026\\[-1pt]{\scriptsize [$+$0.019, $+$0.034]}} & \shortstack{$+$0.028\\[-1pt]{\scriptsize [$+$0.016, $+$0.040]}} & \shortstack{$+$0.027\\[-1pt]{\scriptsize [$+$0.020, $+$0.034]}} & \shortstack{$+$0.028\\[-1pt]{\scriptsize [$+$0.015, $+$0.041]}} \\
Alps & \shortstack{$+$0.026\\[-1pt]{\scriptsize [$+$0.021, $+$0.030]}} & \shortstack{$+$0.029\\[-1pt]{\scriptsize [$+$0.020, $+$0.038]}} & \shortstack{$+$0.026\\[-1pt]{\scriptsize [$+$0.021, $+$0.030]}} & \shortstack{$+$0.032\\[-1pt]{\scriptsize [$+$0.024, $+$0.040]}} \\
\bottomrule
\end{tabular}
\end{adjustbox}
\caption{Difference of $H$, para $-$ hang, in the four designs: fixed population against all flights, over the same lags (to 1000 s, and to the longest $T$ of the population).}\label{tab:fixed-vs-all}
\end{table}


\begin{figure}
\centering
\includegraphics[width=0.62\textwidth]{forest.pdf}
\caption{The paragliders have the larger $H$ in every population and every design, with $\Hpara-\Hhang$ between \UniForestMin{} and \UniForestMax; the band
excludes 0 in \UniForestPos{} of the \UniForestN{} cases. $\Hpara-\Hhang$ with its 90\% band; the number of hang glider flights in
brackets; filled: fixed population (A); open: all flights (B); circle: lags to 1000 s; square: to the longest $T$ of the
population.}\label{fig:forest}
\end{figure}

\keep{The shrinking population lowers $H$ by up to \UniMaxDropB{} and bends $\alpha(\tau)$ at long lags; the
difference between the disciplines is the same in the two designs up to 1000 s, and within the bands to the longest $T$.}
